\section*{Tóm tắt}
\addcontentsline{toc}{section}{Tóm tắt}

Báo cáo này sử dụng bộ dữ liệu World Development Indicators (WDI) của World Bank để khảo sát mức độ phát triển của các quốc gia và khu vực trong giai đoạn 2000--2023. Mục tiêu của báo cáo là xây dựng một khung phân tích chung về phát triển bền vững, trong đó bốn khía cạnh kinh tế, giáo dục, sức khỏe--dân số, và môi trường--biến đổi khí hậu được xem như các nhóm bằng chứng bổ trợ cho nhau. Dữ liệu được tiền xử lý bằng Python, loại bỏ các dòng tổng hợp không đại diện cho quốc gia thực, chuẩn hóa tên biến, và xử lý giá trị thiếu theo quy ước thống nhất trước khi trực quan hóa bằng Tableau.

Nội dung phân tích được cấu trúc mạch lạc thành bốn phần tương ứng với bốn khía cạnh cốt lõi của phát triển bền vững. Phần phân tích kinh tế làm rõ nền tảng nguồn lực vật chất và xu hướng giảm nghèo. Phần giáo dục đánh giá khả năng tiếp cận và mức đầu tư công cho giáo dục. Phần sức khỏe và dân số theo dõi kết quả cải thiện tuổi thọ và giảm tỷ lệ tử vong trẻ em. Cuối cùng, phần môi trường kiểm tra áp lực của tăng trưởng lên phát thải carbon, sự chuyển dịch cơ cấu năng lượng tái tạo và diện tích che phủ rừng. Tất cả các phân tích được tích hợp và trực quan hóa đồng bộ trên một Dashboard Tableau tương tác thống nhất, cho phép người dùng tự do khám phá và so sánh đa chiều.
